\section{Intersecting with an invertible sheaf}
\label{section-intersecting-with-divisors}

\noindent
In this section we study the following construction.

\begin{definition}
\label{definition-cap-c1}
Let $(S, \delta)$ be as in Situation \ref{situation-setup}.
Let $X$ be locally of finite type over $S$.
Let $\mathcal{L}$ be an invertible $\mathcal{O}_X$-module.
We define, for every integer $k$, an operation
$$
c_1(\mathcal{L}) \cap - :
Z_{k + 1}(X) \to \CH_k(X)
$$
called {\it intersection with the first Chern class of $\mathcal{L}$}.
\begin{enumerate}
\item Given an integral closed subscheme $i : W \to X$ with
$\dim_\delta(W) = k + 1$ we define
$$
c_1(\mathcal{L}) \cap [W] = i_*(c_1({i^*\mathcal{L}}) \cap [W])
$$
where the right hand side is defined in
Definition \ref{definition-divisor-invertible-sheaf}.
\item For a general $(k + 1)$-cycle $\alpha = \sum n_i [W_i]$ we set
$$
c_1(\mathcal{L}) \cap \alpha = \sum n_i c_1(\mathcal{L}) \cap [W_i]
$$
\end{enumerate}
\end{definition}

\noindent
Write each $c_1(\mathcal{L}) \cap W_i = \sum_j n_{i, j} [Z_{i, j}]$
with $\{Z_{i, j}\}_j$ a locally finite sum
of integral closed subschemes of $W_i$. Since $\{W_i\}$ is a locally
finite collection of integral closed subschemes on $X$, it follows
easily that $\{Z_{i, j}\}_{i, j}$ is a locally finite collection
of closed subschemes of $X$. Hence
$c_1(\mathcal{L}) \cap \alpha = \sum n_in_{i, j}[Z_{i, j}]$
is a cycle. Another, more convenient, way to think about this
is to observe that the morphism $\coprod W_i \to X$ is
proper. Hence $c_1(\mathcal{L}) \cap \alpha$ can be viewed
as the pushforward of a class in $\CH_k(\coprod W_i) = \prod \CH_k(W_i)$.
This also explains why the result is well defined up to rational
equivalence on $X$.

\medskip\noindent
The main goal for the next few sections is to show that intersecting with
$c_1(\mathcal{L})$ factors through rational equivalence.
This is not a triviality.

\begin{lemma}
\label{lemma-c1-cap-additive}
Let $(S, \delta)$ be as in Situation \ref{situation-setup}.
Let $X$ be locally of finite type over $S$.
Let $\mathcal{L}$, $\mathcal{N}$ be an invertible sheaves on $X$.
Then
$$
c_1(\mathcal{L}) \cap \alpha  + c_1(\mathcal{N}) \cap \alpha =
c_1(\mathcal{L} \otimes_{\mathcal{O}_X} \mathcal{N}) \cap \alpha
$$
in $\CH_k(X)$ for every $\alpha \in Z_{k + 1}(X)$. Moreover,
$c_1(\mathcal{O}_X) \cap \alpha = 0$ for all $\alpha$.
\end{lemma}

\begin{proof}
The additivity follows directly from
Divisors, Lemma \ref{divisors-lemma-c1-additive}
and the definitions. To see that $c_1(\mathcal{O}_X) \cap \alpha = 0$
consider the section $1 \in \Gamma(X, \mathcal{O}_X)$. This restricts
to an everywhere nonzero section on any integral closed subscheme
$W \subset X$. Hence $c_1(\mathcal{O}_X) \cap [W] = 0$ as desired.
\end{proof}

\noindent
Recall that $Z(s) \subset X$ denotes the zero scheme of a global section
$s$ of an invertible sheaf on a scheme $X$, see
Divisors, Definition \ref{divisors-definition-zero-scheme-s}.

\begin{lemma}
\label{lemma-prepare-geometric-cap}
Let $(S, \delta)$ be as in Situation \ref{situation-setup}.
Let $Y$ be locally of finite type over $S$.
Let $\mathcal{L}$ be an invertible $\mathcal{O}_Y$-module.
Let $s \in \Gamma(Y, \mathcal{L})$.
Assume
\begin{enumerate}
\item $\dim_\delta(Y) \leq k + 1$,
\item $\dim_\delta(Z(s)) \leq k$, and
\item for every generic point $\xi$ of an irreducible component of
$Z(s)$ of $\delta$-dimension $k$ the multiplication by $s$
induces an injection $\mathcal{O}_{Y, \xi} \to \mathcal{L}_\xi$.
\end{enumerate}
Write $[Y]_{k + 1} = \sum n_i[Y_i]$ where $Y_i \subset Y$ are the
irreducible components of $Y$ of $\delta$-dimension $k + 1$.
Set $s_i = s|_{Y_i} \in \Gamma(Y_i, \mathcal{L}|_{Y_i})$. Then
\begin{equation}
\label{equation-equal-as-cycles}
[Z(s)]_k =  \sum n_i[Z(s_i)]_k
\end{equation}
as $k$-cycles on $Y$.
\end{lemma}

\begin{proof}
Let $Z \subset Y$ be an integral closed subscheme of
$\delta$-dimension $k$. Let $\xi \in Z$ be its generic point.
We want to compare the coefficient $n$ of $[Z]$ in the expression
$\sum n_i[Z(s_i)]_k$ with the coefficient $m$ of $[Z]$ in the
expression $[Z(s)]_k$. Choose a generator $s_\xi \in \mathcal{L}_\xi$.
Write $A = \mathcal{O}_{Y, \xi}$, $L = \mathcal{L}_\xi$.
Then $L = As_\xi$. Write $s = f s_\xi$ for some (unique) $f \in A$.
Hypothesis (3) means that $f : A \to A$ is injective.
Since $\dim_\delta(Y) \leq k + 1$ and $\dim_\delta(Z) = k$
we have $\dim(A) = 0$ or $1$. We have
$$
m = \text{length}_A(A/(f))
$$
which is finite in either case.

\medskip\noindent
If $\dim(A) = 0$, then $f : A \to A$ being injective
implies that $f \in A^*$. Hence in this case $m$ is zero.
Moreover, the condition $\dim(A) = 0$ means that $\xi$
does not lie on any irreducible component of $\delta$-dimension
$k + 1$, i.e., $n = 0$ as well.

\medskip\noindent
Now, let $\dim(A) = 1$.
Since $A$ is a Noetherian local ring it has finitely
many minimal primes $\mathfrak q_1, \ldots, \mathfrak q_t$.
These correspond 1-1 with the $Y_i$ passing through $\xi'$.
Moreover $n_i = \text{length}_{A_{\mathfrak q_i}}(A_{\mathfrak q_i})$.
Also, the multiplicity of $[Z]$ in $[Z(s_i)]_k$ is
$\text{length}_A(A/(f, \mathfrak q_i))$.
Hence the equation to prove in this case is
$$
\text{length}_A(A/(f))
=
\sum \text{length}_{A_{\mathfrak q_i}}(A_{\mathfrak q_i})
\text{length}_A(A/(f, \mathfrak q_i))
$$
which follows from
Lemma \ref{lemma-additivity-divisors-restricted}.
\end{proof}

\noindent
The following lemma is a useful result in order to compute the intersection
product of the $c_1$ of an invertible sheaf and the cycle associated
to a closed subscheme.
Recall that $Z(s) \subset X$ denotes the zero scheme of a global section
$s$ of an invertible sheaf on a scheme $X$, see
Divisors, Definition \ref{divisors-definition-zero-scheme-s}.

\begin{lemma}
\label{lemma-geometric-cap}
Let $(S, \delta)$ be as in Situation \ref{situation-setup}.
Let $X$ be locally of finite type over $S$.
Let $\mathcal{L}$ be an invertible $\mathcal{O}_X$-module.
Let $Y \subset X$ be a closed subscheme.
Let $s \in \Gamma(Y, \mathcal{L}|_Y)$.
Assume
\begin{enumerate}
\item $\dim_\delta(Y) \leq k + 1$,
\item $\dim_\delta(Z(s)) \leq k$, and
\item for every generic point $\xi$ of an irreducible component of
$Z(s)$ of $\delta$-dimension $k$ the multiplication by $s$
induces an injection
$\mathcal{O}_{Y, \xi} \to (\mathcal{L}|_Y)_\xi$\footnote{For example,
this holds if $s$ is a regular section of $\mathcal{L}|_Y$.}.
\end{enumerate}
Then
$$
c_1(\mathcal{L}) \cap [Y]_{k + 1} = [Z(s)]_k
$$
in $\CH_k(X)$.
\end{lemma}

\begin{proof}
Write
$$
[Y]_{k + 1} = \sum n_i[Y_i]
$$
where $Y_i \subset Y$ are the irreducible components of
$Y$ of $\delta$-dimension $k + 1$ and $n_i > 0$.
By assumption the restriction
$s_i = s|_{Y_i} \in \Gamma(Y_i, \mathcal{L}|_{Y_i})$ is not
zero, and hence is a regular section. By Lemma \ref{lemma-compute-c1}
we see that $[Z(s_i)]_k$ represents $c_1(\mathcal{L}|_{Y_i})$.
Hence by definition
$$
c_1(\mathcal{L}) \cap [Y]_{k + 1} = \sum n_i[Z(s_i)]_k
$$
Thus the result follows from Lemma \ref{lemma-prepare-geometric-cap}.
\end{proof}




